\chapter{Órfãos --- Administração de BD e Temas Coringa}
\label{cap:orfaos-teoria}

Este é o capítulo mais heterogêneo do livro: reúne administração física de
banco de dados (o lado operacional que o Capítulo \ref{cap:banco-dados-teoria}
não cobre, focado que está em modelagem e SQL), arquitetura de computadores
e sistema operacional, sistemas de informação corporativos, e aprendizado de
máquina. São temas que aparecem em prova mas não têm um capítulo natural no
edital do Perfil 3 --- daí o nome. No fim do capítulo há um mapa indicando
onde os demais temas "coringa" (blockchain, virtualização, 2PC etc.) estão
tratados em profundidade nos outros capítulos deste livro.

\section{Administração de Dados $\times$ Administração de Banco de Dados}

\begin{conceito}[Dois papéis com o mesmo objeto e focos opostos]
Os dois cuidam do mesmo recurso --- o dado da organização --- mas em níveis
de abstração diferentes, e é essa diferença de \textbf{nível}, não de
"quem é mais importante", que a prova cobra:

\begin{center}
\begin{tabular}{@{}p{2.2cm}p{4.6cm}p{4.6cm}@{}}
\toprule
& \textbf{Administração de Dados (AD)} & \textbf{Administração de BD (DBA)} \\
\midrule
Foco & \textbf{negócio/lógico} --- o dado como recurso corporativo,
independente de onde ele está armazenado & \textbf{técnico/físico} --- o
SGBD específico funcionando \\
Atividades & modelagem conceitual, políticas de dados, governança,
dicionário de dados & instalação, ajuste de desempenho
(\textit{tuning}), backup, segurança de acesso, disponibilidade \\
Natureza & estratégica, \textbf{independe} do produto de SGBD escolhido &
operacional, \textbf{ligada} a um produto específico (Oracle, PostgreSQL,
SQL Server...) \\
\bottomrule
\end{tabular}
\end{center}

Uma forma de fixar a distinção: a AD decidiria \textit{que} dado a empresa
precisa guardar sobre um cliente e \textit{que} regras de qualidade esse
dado deve seguir; o DBA decidiria \textit{em que servidor}, com
\textit{que configuração de índice} e \textit{que rotina de backup} esse
dado efetivamente vai residir.
\end{conceito}

\begin{cuidado}[Não inverta os dois papéis]
Descrever a AD como "quem cuida do servidor de banco" ou o DBA como "quem
define a política corporativa de qualidade de dados" troca os dois níveis
entre si. A pergunta de desempate: a atividade descrita \textbf{depende de
qual produto de SGBD está instalado}? Se sim, é DBA. Se a resposta seria a
mesma independentemente do SGBD escolhido, é AD.
\end{cuidado}

\section{Recuperação e transações: o protocolo ARIES}

\begin{conceito}[Write-Ahead Logging --- registrar antes de gravar]
A maioria dos SGBDs usa o protocolo \textbf{ARIES} para se recuperar de uma
queda (falta de energia, crash do processo) sem perder transações
confirmadas nem deixar o banco num estado inconsistente. O princípio-base é
o \textbf{WAL} (\textit{Write-Ahead Logging}): antes de gravar uma
alteração \textbf{nos dados em si}, o SGBD grava um registro dessa
alteração \textbf{no log}, em disco. Isso garante a \textbf{durabilidade}
(o D do ACID, ver Capítulo \ref{cap:banco-dados-teoria}) sem precisar
gravar cada página de dado imediatamente a cada operação --- o log, mais
compacto e sequencial, é o que vai para disco com prioridade; o dado em si
pode ser escrito depois, em lote, com mais eficiência.
\end{conceito}

\begin{conceito}[As três fases da recuperação, em ordem fixa]
Depois de uma queda, o ARIES reconstrói o estado do banco em três fases,
\textbf{sempre nesta ordem}:
\begin{enumerate}
\item \textbf{Analysis} (Análise): varre o log a partir do último
\textit{checkpoint} para reconstruir \textbf{que transações} estavam
ativas e \textbf{que páginas} estavam sujas (modificadas, mas não
gravadas) no momento da queda.
\item \textbf{Redo} (Refazer): reaplica \textbf{todas} as operações
registradas no log a partir do ponto identificado --- inclusive as de
transações que \textbf{não chegaram a confirmar} (\textit{commit}). O
princípio é "\textit{repeating history}": primeiro reconstrói o estado
exato de antes da queda, sem julgar ainda o que deveria ou não ter
acontecido.
\item \textbf{Undo} (Desfazer): só depois de reconstruído o estado exato do
momento da queda, desfaz as operações das transações que \textbf{não
tinham confirmado} --- porque uma transação sem commit não deve deixar
efeito permanente.
\end{enumerate}
\end{conceito}

\begin{exemplo}[Por que Redo vem antes de Undo]
Se o Undo rodasse antes do Redo, o sistema tentaria desfazer operações de
transações \textbf{sem primeiro saber com certeza} quais páginas realmente
tinham sido alteradas em disco no momento da queda (algumas mudanças
podem ter chegado ao disco, outras não) --- o Redo existe exatamente para
eliminar essa incerteza, recriando um estado \textbf{conhecido e completo}
antes de começar a reverter o que não deveria ter ficado. Inverter a ordem
quebraria a garantia de consistência que o protocolo promete.
\end{exemplo}

\begin{cuidado}[A troca mais comum: inverter Redo e Undo]
Um enunciado que descreve a recuperação como "primeiro desfaz as
transações não confirmadas, depois refaz as confirmadas" inverteu a lógica
do ARIES. A ordem correta refaz \textbf{tudo primeiro} (inclusive o que
depois será desfeito) e só desfaz \textbf{depois}, num segundo momento.
\end{cuidado}

\section{Armazenamento físico e desempenho}

\begin{conceito}[Tablespace: a unidade lógica que organiza os arquivos físicos]
\textbf{Tablespace} (termo Oracle; \textit{filegroup} no SQL Server) é uma
unidade lógica de armazenamento que agrupa objetos do banco (tabelas,
índices) e os associa a um ou mais arquivos físicos em disco. O motivo de
separar tablespaces distintas para dados, índices e área temporária é
\textbf{isolar o padrão de I/O}: dados e índices costumam ser lidos juntos
mas escritos em ritmos diferentes, e a área temporária (usada por
ordenações e junções grandes) tem um padrão de uso totalmente diferente
--- separar fisicamente evita que um tipo de operação compita com outro
pelo mesmo disco.
\end{conceito}

\begin{conceito}[Otimizador de consultas e o papel das estatísticas]
O \textbf{otimizador de consultas} (\textit{query optimizer}) decide o
\textbf{plano de execução} de um comando SQL --- em que ordem fazer os
\texttt{JOIN}s, se usar um índice ou percorrer a tabela inteira
(\textit{full scan}), qual algoritmo de junção aplicar. Essa decisão se
baseia em \textbf{estatísticas} sobre os dados (quantas linhas tem cada
tabela, quão distribuídos estão os valores de cada coluna) --- não em
executar a consulta de verdade e comparar. Quando as estatísticas ficam
\textbf{desatualizadas} (a tabela cresceu, a distribuição de valores mudou,
mas o otimizador ainda raciocina com o retrato antigo), o plano escolhido
deixa de ser o melhor, e a consulta que era rápida vira lenta sem que
nenhum código tenha mudado. Por isso, \textbf{atualizar as estatísticas} é
frequentemente o primeiro passo de diagnóstico quando uma consulta piora de
repente --- antes de sair reescrevendo SQL ou criando índice novo.
\end{conceito}

\begin{conceito}[Particionamento: dividir para escalar]
\textbf{Particionamento} divide uma tabela grande em pedaços menores e mais
administráveis:
\begin{itemize}
\item \textbf{Horizontal} (também chamado \textit{sharding} quando os
pedaços vão para servidores diferentes): divide por \textbf{linhas} --- por
exemplo, uma tabela de pedidos particionada por ano, cada partição com um
subconjunto das linhas.
\item \textbf{Vertical}: divide por \textbf{colunas} --- separa colunas
acessadas com frequência das raramente acessadas (ou muito grandes, como
um campo de texto longo), para que consultas comuns leiam menos dado.
\end{itemize}
O ganho em ambos os casos é permitir que uma consulta toque só a parte
relevante dos dados, em vez de varrer a tabela inteira.
\end{conceito}

\begin{regrapratica}[Índices: reveja no capítulo de banco de dados]
A escolha entre índice B+ Tree, hash e bitmap já está detalhada na Seção de
Índices do Capítulo \ref{cap:banco-dados-teoria} --- as duas perguntas que
decidem o tipo certo são \textit{a consulta pede faixa ou igualdade?} e
\textit{quantos valores distintos tem a coluna?}. Não repita esse
raciocínio do zero aqui; ele é o mesmo independentemente de o cenário vir
rotulado como "administração de BD" ou "modelagem".
\end{regrapratica}

\section{Backup e recuperação}

\begin{conceito}[Completo, incremental e diferencial --- a diferença é a referência]
Os três tipos de backup diferem em \textbf{a partir de que ponto} eles
medem "o que mudou":

\begin{center}
\begin{tabular}{@{}p{2.6cm}p{8.4cm}@{}}
\toprule
\textbf{Tipo} & \textbf{O que copia} \\
\midrule
\textbf{Completo (full)} & \textbf{tudo}, sem depender de nenhum backup
anterior \\
\textbf{Incremental} & só o que mudou \textbf{desde o último backup, de
qualquer tipo} (o completo mais recente \textbf{ou} o incremental mais
recente) \\
\textbf{Diferencial} & tudo que mudou \textbf{desde o último backup
completo} --- ignora incrementais/diferenciais intermediários \\
\bottomrule
\end{tabular}
\end{center}

A consequência prática é o que costuma ser cobrado: o \textbf{diferencial}
\textbf{cresce} a cada dia que passa sem um novo completo (porque sempre
conta a partir do mesmo ponto fixo), enquanto o \textbf{incremental} tende
a ficar pequeno e constante (porque cada um só cobre o intervalo desde o
anterior, que é sempre recente). Em compensação, restaurar a partir de
incrementais exige reaplicar \textbf{toda a cadeia} de incrementais em
ordem desde o último completo; restaurar a partir de um diferencial exige
só o completo + o diferencial mais recente.
\end{conceito}

\begin{conceito}[RPO e RTO --- as duas perguntas de continuidade de negócio]
\textbf{RPO} (\textit{Recovery Point Objective}): quanto \textbf{dado} a
organização aceita perder, medido em tempo --- se o RPO é de 1 hora, o
backup/replicação precisa garantir que, na pior hipótese, no máximo 1 hora
de transações seja perdida. \textbf{RTO} (\textit{Recovery Time
Objective}): em quanto \textbf{tempo} o serviço precisa voltar a operar
depois de um incidente. São eixos independentes: um sistema pode ter RPO
baixo (não pode perder quase nada de dado) mas RTO alto (pode ficar fora do
ar por horas), ou o contrário.
\end{conceito}

\begin{conceito}[Backup quente $\times$ frio]
\textbf{Quente} (\textit{online}): feito com o banco \textbf{no ar},
atendendo requisições normalmente durante a cópia. \textbf{Frio}
(\textit{offline}): feito com o banco \textbf{parado} --- mais simples de
garantir consistência, mas exige uma janela de indisponibilidade.
\end{conceito}

\section{Segurança e acesso no SGBD}

\begin{conceito}[Camadas de controle de acesso dentro do banco]
\textbf{GRANT/REVOKE} (comandos DCL, ver Capítulo
\ref{cap:banco-dados-teoria}) concedem e revogam privilégios diretamente.
\textbf{Papéis/roles} (RBAC --- \textit{Role-Based Access Control}) agrupam
privilégios e são atribuídos a usuários, evitando conceder/revogar
privilégio individual a cada pessoa. \textbf{Views} podem restringir o
acesso a \textbf{colunas} específicas (o usuário enxerga só um subconjunto
de campos de uma tabela). \textit{\textbf{Row-level security}} restringe o
acesso por \textbf{linha} (ex.: o gerente de uma filial só vê as linhas
daquela filial). Auditoria registra quem acessou o quê, para rastreamento
posterior.

O princípio que organiza tudo isso é o \textbf{menor privilégio}: cada
usuário/aplicação recebe exatamente o acesso necessário para sua função,
nada além --- o gerente de RH acessa dados de folha de pagamento, não
acessa esquema de outro sistema que não lhe diz respeito.
\end{conceito}

\section{Arquitetura de computadores e sistema operacional}

\begin{conceito}[Fora do escopo estrito do edital, mas cai em prova]
O edital do Perfil 3 fala em arquitetura \emph{de software}
(Capítulo \ref{cap:arquitetura-teoria}), não em arquitetura \emph{de
hardware}. Ainda assim, provas de concurso costumam reaproveitar questões
de "informática básica"/arquitetura de computadores entre perfis
diferentes, e esse conteúdo aparece. É rápido de cobrir e vale a pena
saber.
\end{conceito}

\subsection{Modelo de Von Neumann}

\begin{conceito}[Um só barramento para dados e instruções --- e o gargalo que isso cria]
A arquitetura de Von Neumann descreve um computador com \textbf{memória
única} compartilhada entre \textbf{instruções} (o programa) e
\textbf{dados} (as informações que o programa manipula), acessada por um
\textbf{único barramento}. A CPU busca uma instrução, depois busca/grava
dados, pelo mesmo canal --- nunca simultaneamente. Esse compartilhamento é
conhecido como o \textbf{gargalo de Von Neumann}: por mais rápida que a CPU
fique, ela é limitada pela velocidade com que consegue trafegar instrução e
dado pelo mesmo barramento, um de cada vez.
\end{conceito}

\subsection{CPU: Unidade de Controle e ULA}

\begin{conceito}[Duas unidades, duas responsabilidades]
\begin{itemize}
\item \textbf{Unidade de Controle (UC)}: \textbf{coordena} a execução ---
busca a instrução, decodifica o que ela significa, e ativa os sinais de
controle necessários para executá-la (não realiza o cálculo em si, apenas
orquestra).
\item \textbf{Unidade Lógica e Aritmética (ULA)}: executa as operações
\textbf{de fato} --- soma, subtração, comparações lógicas (E, OU, NÃO).
\end{itemize}
A analogia útil: a UC é o maestro (decide o que tocar e quando), a ULA é
quem toca o instrumento.
\end{conceito}

\subsection{Ciclo de instrução e \textit{pipeline}}

\begin{conceito}[Busca, decodifica, executa]
Toda instrução de máquina passa por um ciclo básico:
\textbf{busca} (\textit{fetch} --- traz a instrução da memória),
\textbf{decodifica} (\textit{decode} --- interpreta o que ela pede) e
\textbf{executa} (\textit{execute} --- realiza a operação, geralmente na
ULA). Um processador simples faz esse ciclo inteiro para uma instrução
antes de começar o próximo.
\end{conceito}

\begin{conceito}[Pipeline: sobrepor os ciclos]
Um processador com \textbf{pipeline} não espera uma instrução terminar
todo o ciclo antes de começar a próxima: enquanto a instrução A está sendo
\textbf{executada}, a instrução B já está sendo \textbf{decodificada} e a
instrução C já está sendo \textbf{buscada} --- como uma linha de montagem,
em que cada estação trabalha numa instrução diferente ao mesmo tempo. O
ganho é em \textbf{taxa de conclusão de instruções por segundo}
(\textit{throughput}), não em fazer uma instrução isolada terminar mais
rápido.
\end{conceito}

\subsection{Sistemas de numeração e complemento de dois}

\begin{conceito}[Conversão entre bases --- o mesmo valor, notações diferentes]
Um computador só entende binário (base 2) internamente, mas humanos leem
mais facilmente octal (base 8) ou hexadecimal (base 16) --- ambas
escolhidas por serem \textbf{potências de 2} ($8=2^3$, $16=2^4$), o que
torna a conversão \textbf{dígito a dígito} sem precisar de conta longa:
agrupe os bits binários em grupos de 3 (octal) ou 4 (hexadecimal), da
direita para a esquerda, e converta cada grupo isoladamente.
\end{conceito}

\begin{exemplo}[Convertendo binário para hexadecimal]
Binário: \texttt{1011 1100}. Agrupe de 4 em 4, da direita para a esquerda:
\texttt{1011} e \texttt{1100}. \texttt{1011} = 8+0+2+1 = 11 = \texttt{B}.
\texttt{1100} = 8+4+0+0 = 12 = \texttt{C}. Resultado: \texttt{0xBC}.
\end{exemplo}

\begin{conceito}[Complemento de dois: como o computador representa negativos]
Para representar números negativos em binário sem precisar de um circuito
separado para subtração, usa-se o \textbf{complemento de dois}: para negar
um número, inverte-se todos os bits (complemento de um) e soma-se 1. A
vantagem é que \textbf{a mesma unidade de soma} da ULA resolve subtração
também --- "$A - B$" vira "$A + (\text{complemento de dois de } B)$", sem
circuito dedicado.
\end{conceito}

\begin{conceito}[\textit{Overflow} $\times$ \textit{carry} --- dois tipos de "estourou"]
\begin{itemize}
\item \textbf{Carry} (vai-um): ocorre quando uma soma \textbf{sem sinal}
ultrapassa a capacidade de bits disponível --- é o "vai um" que sobra da
última posição.
\item \textbf{Overflow}: ocorre quando o resultado de uma operação
\textbf{com sinal} (complemento de dois) produz um valor que \textbf{não
cabe} na faixa representável com aquele número de bits --- por exemplo,
somar dois números positivos grandes e o resultado "virar" negativo por
estourar o bit de sinal.
\end{itemize}
Os dois podem ocorrer independentemente um do outro na mesma operação ---
não são sinônimos, e um enunciado que os trata como a mesma coisa está
simplificando demais um conceito que a prova cobra separado.
\end{conceito}

\subsection{ROM, firmware e periféricos}

\begin{conceito}[Memória que sobrevive sem energia]
\textbf{ROM} (\textit{Read-Only Memory}) é memória \textbf{não volátil}
(mantém o conteúdo sem energia), tipicamente usada para armazenar
\textbf{firmware} --- o software de baixo nível gravado de fábrica que
inicializa o hardware antes mesmo do sistema operacional carregar (ex.: a
BIOS/UEFI de um computador). \textbf{Periférico} é qualquer dispositivo de
entrada/saída conectado à CPU (teclado, monitor, disco), que se comunica
com o resto do sistema através de controladores e barramentos de E/S.
\end{conceito}

\subsection{Memória virtual: paginação, MMU, TLB e \textit{thrashing}}

\begin{conceito}[O problema: programas maiores que a RAM física disponível]
\textbf{Memória virtual} é a técnica que permite a um programa "enxergar"
mais memória do que a RAM física realmente tem, usando o disco como
extensão. O sistema operacional divide a memória em blocos de tamanho fixo
chamados \textbf{páginas}, e mantém uma \textbf{tabela de páginas} que
mapeia \textbf{endereço virtual} (o que o programa usa) para
\textbf{endereço físico} (onde o dado realmente está --- na RAM ou, se não
estiver presente no momento, no disco).
\end{conceito}

\begin{conceito}[MMU e TLB: o hardware que acelera a tradução]
A \textbf{MMU} (\textit{Memory Management Unit}) é o circuito de hardware
que traduz endereço virtual em físico a cada acesso à memória, consultando
a tabela de páginas. Como consultar a tabela de páginas na RAM a
\textbf{cada} acesso seria lento, a MMU mantém um pequeno cache dedicado
chamado \textbf{TLB} (\textit{Translation Lookaside Buffer}) com as
traduções mais recentes/frequentes --- um acerto na TLB evita ter que
consultar a tabela de páginas completa.
\end{conceito}

\begin{conceito}[\textit{Thrashing}: quando a paginação vira o próprio problema]
Se a memória física disponível é pequena demais para as páginas que os
processos ativos realmente precisam a cada instante, o sistema passa a
gastar mais tempo \textbf{trocando páginas entre RAM e disco}
(\textit{swapping}) do que executando trabalho útil --- esse colapso de
desempenho é o \textbf{thrashing}. O sintoma característico é a CPU parecer
\textbf{ociosa} mesmo com processos "rodando": na verdade eles estão
travados esperando página chegar do disco, não executando instrução.
\end{conceito}

\subsection{E/S bloqueante, troca de contexto e DMA}

\begin{conceito}[E/S bloqueante x troca de contexto]
Quando um processo pede uma operação de \textbf{entrada/saída bloqueante}
(ex.: ler um arquivo em disco), ele fica \textbf{suspenso} esperando a
operação terminar --- não faz sentido a CPU ficar parada nesse meio tempo,
então o sistema operacional aproveita para dar a CPU a \textbf{outro}
processo pronto para rodar. Essa troca --- salvar o estado do processo
suspenso e carregar o estado do próximo --- é a \textbf{troca de contexto}
(\textit{context switch}), o mecanismo que viabiliza multitarefa numa CPU
que só executa uma instrução por vez em cada núcleo.
\end{conceito}

\begin{conceito}[DMA: tirar a CPU do meio do caminho]
\textbf{DMA} (\textit{Direct Memory Access}) é um controlador dedicado que
transfere dados diretamente entre um periférico (ex.: disco) e a memória
RAM, \textbf{sem} que a CPU precise mediar byte a byte a transferência ---
a CPU apenas configura a operação (origem, destino, tamanho) e é
\textbf{interrompida} quando a transferência termina, ficando livre para
fazer outra coisa enquanto o DMA trabalha. Sem DMA, toda transferência de
E/S consumiria ciclos de CPU copiando dado por dado.
\end{conceito}

\section{Sistemas de informação corporativos}

\begin{conceito}[Cada sistema atende um nível gerencial diferente]
Sistemas de informação corporativos se diferenciam pelo \textbf{nível de
decisão} que atendem --- do operacional ao estratégico:
\begin{itemize}
\item \textbf{SPT} (Sistema de Processamento de Transações, também
chamado transacional): registra as operações do \textbf{dia a dia}
operacional (uma venda, um pagamento) --- nível mais baixo, alto volume,
baixa complexidade de decisão por transação.
\item \textbf{SIG} (Sistema de Informação Gerencial): consolida dados
operacionais em \textbf{relatórios} para a gerência tática acompanhar
desempenho.
\item \textbf{SAD/SSD} (Sistema de Apoio/Suporte à Decisão): fornece
\textbf{análise} para decisões \textbf{não estruturadas} ou semiestruturadas,
tipicamente de nível mais estratégico --- simulações, cenários "e se".
\item \textbf{ERP} (\textit{Enterprise Resource Planning}): integra
processos de \textbf{múltiplas} áreas (financeiro, estoque, produção,
RH) num único sistema com base de dados compartilhada, eliminando
retrabalho de dados duplicados entre sistemas isolados.
\item \textbf{CRM} (\textit{Customer Relationship Management}): foca no
relacionamento com o \textbf{cliente} --- histórico de contato, vendas,
suporte.
\item \textbf{SCM} (\textit{Supply Chain Management}): foca na
\textbf{cadeia de suprimentos} --- fornecedores, logística, estoque entre
elos da cadeia.
\end{itemize}
A pirâmide organizacional clássica (operacional na base, tático no meio,
estratégico no topo) ajuda a lembrar: SPT nasce na base, SIG no meio, SAD
no topo; ERP/CRM/SCM são \textbf{transversais}, cortando vários níveis ao
mesmo tempo, cada um por um recorte de processo diferente (recursos
internos, cliente, cadeia de fornecimento).
\end{conceito}

\section{Aprendizado de máquina}

\begin{conceito}[Supervisionado, não supervisionado e semissupervisionado]
A distinção central é se os dados de treino trazem \textbf{rótulo}
(\textit{label} --- a resposta certa já conhecida):
\begin{itemize}
\item \textbf{Supervisionado}: treina com dados \textbf{rotulados} ---
cada exemplo já vem com a resposta esperada. Usado em
\textbf{classificação} (rótulo categórico: "é spam?") e \textbf{regressão}
(rótulo numérico: "qual será o preço?").
\item \textbf{Não supervisionado}: treina com dados \textbf{sem
rótulo} --- o algoritmo busca \textbf{estrutura} escondida nos próprios
dados, sem saber de antemão qual é a "resposta certa". Usado em
\textbf{clusterização} (agrupar exemplos parecidos) e \textbf{associação}
(achar padrões de coocorrência).
\item \textbf{Semissupervisionado}: usa uma \textbf{pequena} quantidade de
dados rotulados junto com uma \textbf{grande} quantidade de dados sem
rótulo --- útil quando rotular é caro (exige um humano), mas coletar dado
bruto é barato.
\end{itemize}
\end{conceito}

\begin{cuidado}[Não inverta "rotulado" entre os dois]
Um enunciado que descreve o \textbf{não supervisionado} como "usa dados
rotulados para prever categorias" trocou a definição com o
\textbf{supervisionado}. O critério de desempate é sempre: \textbf{existe
resposta certa conhecida de antemão no conjunto de treino, ou o algoritmo
tem que descobrir a estrutura sozinho?}
\end{cuidado}

\begin{conceito}[Overfitting e underfitting]
\textbf{Overfitting}: o modelo \textbf{decora} particularidades do
conjunto de treino (inclusive ruído) em vez de aprender o padrão geral ---
desempenho ótimo no treino, ruim em dados novos. \textbf{Underfitting}: o
modelo é \textbf{simples demais} para captar sequer o padrão do treino ---
desempenho ruim em ambos, treino e dados novos. O tratamento detalhado da
relação viés$\times$variância por trás desses dois fenômenos está no
Capítulo \ref{cap:atualidades-teoria}.
\end{conceito}

\begin{conceito}[Regularização L1 e L2: como conter o overfitting pela penalidade]
Um modelo \textit{overfit} costuma ter \textbf{coeficientes grandes} --- é
o que permite a curva ajustada oscilar bruscamente para passar exatamente
por cada ponto do treino. Regularizar significa treinar minimizando
\textbf{erro + uma penalidade sobre o tamanho dos coeficientes}, tornando
caro manter coeficiente grande sem necessidade real:

\begin{center}
\begin{tabular}{@{}p{2.4cm}p{3.6cm}p{5.4cm}@{}}
\toprule
& \textbf{Penaliza} & \textbf{Efeito} \\
\midrule
\textbf{L1 (Lasso)} & soma dos \textbf{valores absolutos} dos coeficientes
& empurra coeficientes \textbf{até zero} --- remove atributos do modelo
(\textbf{seleção de atributos}, modelo \textit{esparso}) \\
\textbf{L2 (Ridge)} & soma dos \textbf{quadrados} dos coeficientes &
\textbf{encolhe} todos proporcionalmente, mas \textbf{não zera nenhum} ---
todos os atributos permanecem, com peso menor \\
\bottomrule
\end{tabular}
\end{center}

A razão da diferença está no \textbf{formato} da penalidade: elevar ao
quadrado (L2) pune muito um coeficiente grande e quase nada um já pequeno
--- não compensa terminar de zerá-lo. A L1 pune o mesmo valor por unidade
em qualquer faixa, então zerar de vez sai "de graça" e compensa. Use L1
quando o objetivo é \textbf{descartar} atributos irrelevantes de um
conjunto com muitas variáveis; use L2 quando o objetivo é apenas
\textbf{estabilizar} um modelo com atributos correlacionados, mantendo
todos.
\end{conceito}

\begin{conceito}[Métricas de avaliação]
\textbf{MAE} (\textit{Mean Absolute Error}): erro médio absoluto, para
\textbf{regressão} --- a distância média entre previsão e valor real, sem
distinguir se errou para mais ou para menos. \textbf{F1-score}: média
\textbf{harmônica} entre precisão (dos que o modelo disse positivo, quantos
eram de fato) e recall/revocação (dos que eram de fato positivos, quantos o
modelo capturou) --- usada quando as classes são desbalanceadas e acurácia
simples enganaria. \textbf{ROC/AUC}: mede a capacidade do modelo de
\textbf{separar classes} variando o limiar de decisão --- quanto mais perto
de 1 a área sob a curva (AUC), melhor a separação.
\end{conceito}

\begin{conceito}[Drift: o modelo não piora sozinho, o mundo muda ao redor]
Um modelo em produção perde acurácia ao longo do tempo mesmo sem nenhuma
mudança de código, porque foi treinado numa "fotografia" da realidade que
deixa de valer. Duas causas distintas:

\begin{center}
\begin{tabular}{@{}p{2.6cm}p{4.2cm}p{4.8cm}@{}}
\toprule
& \textbf{O que mudou} & \textbf{Exemplo} \\
\midrule
\textbf{Data drift} & a \textbf{distribuição da entrada} --- os dados que
chegam já não se parecem com os de treino & o perfil de usuários do
sistema mudou (faixa etária, região), mas a relação entre perfil e
comportamento continua a mesma \\
\textbf{Concept drift} & a \textbf{relação entre entrada e saída} --- o
próprio padrão que o modelo aprendeu deixou de valer & uma lei nova (ou
mudança de mercado) faz o mesmo perfil de cliente se comportar de um jeito
diferente do que se comportava antes \\
\bottomrule
\end{tabular}
\end{center}

A diferença tem consequência prática: data drift às vezes se resolve
reamostrando/rebalanceando os dados de treino; concept drift \textbf{exige
retreinar com rótulos novos}, porque a resposta certa em si mudou --- não
adianta só ajustar a proporção dos dados antigos. Detectar os dois é papel
do monitoramento contínuo de modelos em produção (MLOps).
\end{conceito}

\section{Mapa: onde estão os demais temas coringa}

Vários temas que também caem sob a etiqueta "sem capítulo próprio" já
recebem tratamento profundo em outros capítulos deste livro --- não os
repita aqui, use este mapa:

\begin{center}
\begin{tabular}{@{}p{5cm}p{6cm}@{}}
\toprule
\textbf{Tema} & \textbf{Onde está} \\
\midrule
Blockchain (hash encadeado, raiz de Merkle, \textit{nonce}) & Capítulo
\ref{cap:programacao-teoria} \\
Servidor web $\times$ servidor de aplicação & Capítulo
\ref{cap:arquitetura-teoria} \\
2PC (\textit{two-phase commit}) $\times$ Raft (consenso) & Capítulo
\ref{cap:arquitetura-teoria}, Seção~\ref{sec:transacoes-dist} \\
Virtualização: hipervisor Tipo 1 $\times$ Tipo 2 $\times$ contêiner &
Capítulo \ref{cap:arquitetura-teoria} \\
Low-code / no-code & Capítulo \ref{cap:programacao-teoria} \\
RPA (\textit{Robotic Process Automation}) & Capítulo
\ref{cap:eng-software-teoria} \\
Os "V"s do Big Data & Capítulo \ref{cap:banco-dados-teoria} \\
Índices (B+ Tree, hash, bitmap) & Capítulo \ref{cap:banco-dados-teoria} \\
Internet $\times$ intranet $\times$ extranet $\times$ portal & Capítulo
\ref{cap:redes-teoria} \\
IA generativa, LLM, RAG, ética e regulação & Capítulo
\ref{cap:atualidades-teoria} \\
\bottomrule
\end{tabular}
\end{center}

\section{Roteiro de revisão do capítulo}

\begin{regrapratica}[Pares que se invertem com frequência]
Memorize estes pares em bloco, porque a confusão entre os dois lados de
cada par é o erro mais comum deste capítulo: AD (negócio) $\times$ DBA
(técnico); incremental (desde o último backup qualquer) $\times$
diferencial (desde o último completo); RPO (quanto dado perder) $\times$
RTO (quanto tempo até voltar); ordem Analysis $\to$ Redo $\to$ Undo do
ARIES; carry (sem sinal) $\times$ overflow (com sinal); supervisionado
(rotulado) $\times$ não supervisionado (sem rótulo); L1 (zera, seleciona)
$\times$ L2 (encolhe, mantém todos); data drift (mudou quem entra)
$\times$ concept drift (mudou a regra).
\end{regrapratica}
